% Part: proof-theory
% Chapter: sequent-calculus
% Section: invertibility

\documentclass[../../../include/open-logic-section]{subfiles}

\begin{document}

\olfileid{pt}{seq}{inv}

\olsection{Kuwalikan Aturan}

\begin{defn}
Aturan~$R$,
\[
\AxiomC{$S_1$}
\AxiomC{$\dots$}
\AxiomC{$S_n$}
\RightLabel{$R$}
\TrinaryInfC{$S$}
\DisplayProof
\]
diarani \emph{bisa diwalik} ing sistem yen saben $\Proves S$,
uga $\Proves S_i$ kanggo $i = 1$, \dots,~$n$.
Aturan iki \emph{bisa diwalik kanthi njaga !!{height}}
(utawa \emph{bisa diwalik-!!{hp}}) ing sistem yen saben
$\Proves_h S$, uga $\Proves[h] S_i$ kanggo $i = 1$, \dots,~$n$.
\end{defn}

\begin{lem}\ollabel{lem:G3c-invert}
Aturan~\Log{G3c} kanggo $\lnot$, $\land$, $\lor$, lan $\lif$
bisa diwalik kanthi njaga !!{height}, yaiku:
\begin{enumerate}
  \item \RightR{\lnot}: Yen $\Log{G3c} \Proves[n] \Gamma \Sequent
  \Delta, \lnot !A$, mula $\Log{G3c} \Proves[n] !A, \Gamma \Sequent
  \Delta$.
  \item \LeftR{\lnot}: Yen $\Log{G3c} \Proves[n] \Gamma, \lnot !A
  \Sequent \Delta$, mula $\Log{G3c} \Proves[n] \Gamma \Sequent \Delta,
  !A$.
  \item \RightR{\land}: Yen $\Log{G3c} \Proves[n] \Gamma \Sequent
  \Delta, !A \land !B$, mula $\Log{G3c} \Proves[n] \Gamma \Sequent
  \Delta, !A$ lan $\Log{G3c} \Proves[n] \Gamma \Sequent
  \Delta, !B$.
  \item \LeftR{\land}: Yen $\Log{G3c} \Proves[n] \Gamma, !A \land !B
  \Sequent \Delta$, mula $\Log{G3c} \Proves[n] !A, !B, \Gamma \Sequent
  \Delta$.
  \item \RightR{\lor}: Yen $\Log{G3c} \Proves[n] \Gamma \Sequent
  \Delta, !A \lor !B$, mula $\Log{G3c} \Proves[n] \Gamma \Sequent
  \Delta, !A, !B$.
  \item \LeftR{\lor}: Yen $\Log{G3c} \Proves[n] !A \lor !B, \Gamma \Sequent
  \Delta$, mula $\Log{G3c} \Proves[n] !A, \Gamma \Sequent
  \Delta$ lan $\Log{G3c} \Proves[n] !B, \Gamma \Sequent
  \Delta$.
  \item \RightR{\lif}: Yen $\Log{G3c} \Proves[n] \Gamma \Sequent
  \Delta, !A \lif !B$, mula $\Log{G3c} \Proves[n] !A, \Gamma \Sequent
  \Delta, !B$.
  \item \LeftR{\lif}: Yen $\Log{G3c} \Proves[n] !A \lif !B, \Gamma
  \Sequent \Delta$, mula $\Log{G3c} \Proves[n] \Gamma \Sequent \Delta,
  !A$ lan $\Log{G3c} \Proves[n] !B, \Gamma \Sequent \Delta$.
\end{enumerate}
\end{lem}

\begin{proof}
Kita kudu nuduhake yen kanggo saben aturan proposisional~$R$
saka~\Log{G3c}, yen ana !!a{proof}~$\pi$ kanggo kesimpulan~$S$
saka~$R$, mula uga ana !!{proof}s~$\pi_i$ kanggo premis $S_i$
saka~$R$, kanthi $\pheight{\pi_i} \le \pheight{\pi}$.
Gatekna aturan \LeftR{\land}: upamane ana !!a{proof}~$\pi$
kanthi sekuen pungkasan awujud $!A \land !B, \Gamma \Sequent \Delta$.
Kita kudu nuduhake yen ana !!a{proof} $\pi'$ kanggo
$!A, !B, \Gamma \Sequent \Delta$ kanthi
$\pheight{\pi'} \le \pheight{\pi}$. Iki ditindakake kanthi
induksi tumrap $n = \pheight{\pi}$.

Yen $n = 0$, $\pi$ mung kasusun saka aksioma
$!A \land !B, \Gamma \Sequent \Delta$. Ing \Log{G3c},
aksioma identitas awujud $!C, \Gamma \Sequent \Delta, !C$
kanthi $!C$ atomik. Kanthi tembung liya, $!C$ ora bisa
dadi $!A \land !B$. Yen $!A \land !B, \Gamma \Sequent \Delta$
minangka aksioma, mula ana formula atomik $!C$ sing dadi
!!a{element} ing $\Gamma$ lan~$\Delta$, utawa ana falsitas
ing konteks anteseden. Mula $!A, !B, \Gamma \Sequent \Delta$
uga aksioma, lan kita wis nemokake !!{proof}~$\pi'$ sing
uga duwe !!{height}~$0$.

Yen $n>0$, bukti dipungkasi nganggo inferensi. Kita kudu
mbuktekake pratelan kanggo~$n$, nanging bisa nganggep pratelan
iki bener kanggo saben !!{proof} sing dhuwure~$<n$, sanajan
kontekse beda. Kanthi tembung liya, kanggo saben $k<n$ lan
saben $\Pi$ lan $\Lambda$, yen $!A \land !B, \Pi \Sequent \Lambda$
duwe !!a{proof} kanthi !!{height}~$k$, mula ana !!a{proof}
kanthi !!{height}~$\le k$ kanggo $!A, !B, \Pi \Sequent \Lambda$.

Ana rong kasus: $!A \land !B$ ing sisih kiwa dadi formula
utama, utawa ora. Yen dadi formula utama, inferensi pungkasan
yaiku \LeftR{\land}, lan sub-!!{proof} langsung~$\pi_1$
dipungkasi ing $!A, !B, \Gamma \Sequent \Delta$. Ing kasus iki
asile langsung ditampa amarga $\pheight{\pi_1} < \pheight{\pi}$.

Yen $!A \land !B$ dudu formula utama, mula ana !!{formula}
liyane sing dadi formula utama, lan $!A \land !B$ kalebu konteks
inferensi pungkasan. Hipotesis induksi bisa ditrapake marang
sub-!!{proof}s langsung saka aturan pungkasan~$R$, lan menehi
!!a{proof}~$\pi_1'$ utawa rong !!{proof}s~$\pi_1'$ lan~$\pi_2'$
kanthi !!{height} sing ora ngluwihi dhuwur sub-!!{proof}s
langsung saka~$\pi$. Sekuen pungkasan $\pi_1'$ lan $\pi_2'$
ngemot $!A, !B$ ing konteks kiwa minangka gantine $!A \land !B$.
Trapna aturan~$R$ kanggo oleh bukti~$\pi'$. Upamane,
$\pi$ dipungkasi nganggo~\LeftR{\lif}:
\[
\AxiomC{}
\RightLabel{$\pi_1$}
\Deduce$!A \land !B, \Gamma' \fCenter \Delta, !C$
\AxiomC{}
\RightLabel{$\pi_2$}
\Deduce$!A \land !B, \Gamma', !D \fCenter \Delta$
\RightLabel{\LeftR{\lif}}
\BinaryInf$!A \land !B, \Gamma', !C \lif !D \fCenter \Delta$
\DisplayProof
\]
Ing kene $!C \lif !D$ ing kiwa dadi formula utama, lan
konteks kiwa yaiku $!A \land !B, \Gamma'$ (yaiku
$\Gamma = \Gamma', !C \lif !D$). Hipotesis induksi menehi
!!{proof}s $\pi_1'$ lan $\pi_2'$ kanggo
$!A, !B, \Gamma' \Sequent \Delta, !C$ lan
$!A, !B, \Gamma', !D \Sequent \Delta$, manut urutane.
Rong sekuen iki bisa maneh dadi premis aturan \LeftR{\lif}
kanggo oleh~$\pi'$:
\[
\AxiomC{}
\RightLabel{$\pi_1'$}
\Deduce$!A, !B, \Gamma' \fCenter \Delta, !C$
\AxiomC{}
\RightLabel{$\pi_2'$}
\Deduce$!A, !B, \Gamma', !D \fCenter \Delta$
\RightLabel{\LeftR{\lif}}
\BinaryInf$!A, !B, \Gamma', !C \lif !D \fCenter \Delta$
\DisplayProof
\]
Kita duwe
\begin{align*}
\pheight{\pi} & = \max(\pheight{\pi_1}, \pheight{\pi_2}) + 1\\
\pheight{\pi'} & = \max(\pheight{\pi_1'}, \pheight{\pi_2'}) + 1
\intertext{miturut dhefinisi, lan}
\pheight{\pi_1'} & \le \pheight{\pi_1}\\
\pheight{\pi_2'} & \le \pheight{\pi_2}
\intertext{miturut hipotesis induksi. Dadi,}
\pheight{\pi'} & \le \pheight{\pi}.
\end{align*}
\end{proof}

\begin{prob}
Wenehana bukti \olref{lem:G3c-invert} kanggo \RightR{\lif}.
\end{prob}

\begin{lem}\ollabel{lem:invert-quant}
Aturan \RightR{\lforall} lan \LeftR{\lexists} bisa diwalik
kanthi njaga !!{height} ing pangerten iki:
\begin{enumerate}
  \item Yen $\Log{G3c} \Proves[n] \Gamma \Sequent \Delta,
  \lforall[x][!A(x)]$, mula $\Log{G3c} \Proves[n] \Gamma \Sequent
  \Delta, !A(c)$, lan
  \item yen $\Log{G3c} \Proves[n] \lexists[x][!A(x)], \Gamma \Sequent
  \Delta$, mula $\Log{G3c} \Proves[n] !A(c), \Gamma \Sequent \Delta$,
\end{enumerate}
kanggo saben $c$ sing ora ana ing $\Gamma$, $\Delta$, lan
$\lforall[x][!A(x)]$ utawa $\lexists[x][!A(x)]$, manut urutane.
\end{lem}

\begin{proof}
  Kita nuduhake (1) lan ninggalake (2) minangka latihan.
  Kasus dasar aksioma identitas lan falsitas sing padha uga
  ditrapake. Dhisik, ganti jeneng eigenparameter ing subwit
  buktine supaya seger tumrap konstanta sing dipilih.
  Argumentasine padha karo bukti \olref{lem:G3c-invert}.
  Ing langkah induksi, ana rong kasus maneh. Kasus nalika
  $\lforall[x][!A(x)]$ dadi formula utama bisa dirampungake
  langsung: premis inferensi $\RightR{\lforall}$ pungkasan
  wis awujud sing dibutuhake, yaiku $\Gamma \Sequent \Delta, !A(a)$,
  lan sub-!!{proof} langsung~$\pi_1$ sing dipungkasi ing
  premis kasebut duwe $\pheight{\pi_1} < \pheight{\pi}$.
  Kajaba iku, syarat eigenvariabel njamin $a$ ora ana ing
  $\Gamma$, $\Delta$, utawa $\lforall[x][!A(x)]$.
  Amarga kita bisa nganggep $\pi_1$ reguler,
  $\Subst{\pi_1}{c}{a}$ minangka !!a{proof} kanthi
  !!{height} sing padha kanggo $\Gamma \Sequent \Delta, !A(c)$,
  miturut \olref[qua]{lem:subst-regular}.

  Ing kasus kapindho, $\lforall[x][!A(x)]$ dudu formula utama
  ing inferensi pungkasan; formula liyane sing dadi formula
  utama. Kita menehi siji kasus minangka conto maneh.
  Upamane inferensi pungkasan yaiku~\RightR{\land}.
  Ana !!{formula} ing $\Delta$ awujud~$!B \land !C$ sing
  dadi formula utama. !!{proof}~$\pi$ dipungkasi nganggo
  \[
\AxiomC{}
\RightLabel{$\pi_1$}
\Deduce$\Gamma \fCenter \Delta', \lforall[x][!A(x)], !B$
\AxiomC{}
\RightLabel{$\pi_2$}
\Deduce$\Gamma \fCenter \Delta', \lforall[x][!A(x)], !C$
\RightLabel{\RightR{\land}}
\BinaryInf$\Gamma \fCenter \Delta', \lforall[x][!A(x)], !B \land !C$
\DisplayProof
\]
kanthi $\Delta = \Delta', !B \land !C$. Pilih $c$ minangka
!!a{constant} sing nyukupi kabeh syarat kesegaran ing ndhuwur;
mligine, konstanta iki ora ana ing $\Delta$. Dadi cetha yen
uga ora ana ing $\Delta', !B$ utawa ing $\Delta', !C$.
Mula hipotesis induksi bisa ditrapake marang $\pi_1$ lan~$\pi_2$;
kita duwe !!{proof}s~$\pi_1'$ lan $\pi_2'$ kanthi
$\pheight{\pi_1'} \le \pheight{\pi_1}$ lan
$\pheight{\pi_2'} \le \pheight{\pi_2}$. !!{proof}~$\pi'$
sing dibutuhake kanggo $\Gamma \Sequent \Delta, !A(c)$ yaiku:
\[
\AxiomC{}
\RightLabel{$\pi_1'$}
\Deduce$\Gamma \fCenter \Delta', !A(c), !B$
\AxiomC{}
\RightLabel{$\pi_2'$}
\Deduce$\Gamma \fCenter \Delta', !A(c), !C$
\RightLabel{\RightR{\land}}
\BinaryInf$\Gamma \fCenter \Delta', !A(c), !B \land !C$
\DisplayProof
\]
kanthi $\pheight{\pi'} = \max(\pheight{\pi_1'}, \pheight{\pi_2'})+1 \le \pheight{\pi}$.
\end{proof}

\Olref{lem:G3c-invert} duwe peran wigati ing bukti teorema
eliminasi potongan. Nanging lemma iki uga bisa digunakake
kanggo nuduhake asil iki:

\begin{prop}\ollabel{prop:G3c-cont-adm}
Aturan kontraksi \LeftR{\Contraction} lan \RightR{\Contraction}
saka~\Log{G1c} admisibel kanthi njaga !!{height} ing~\Log{G3c}.
\end{prop}

\begin{proof}
Kita menehi argumentasi kanggo \RightR{\Contraction} lan
\LeftR{\Contraction} bebarengan. Upamane $\pi$ minangka
!!a{proof} kanggo $\Gamma \Sequent \Delta, !A, !A$ utawa
$!A, !A, \Gamma \Sequent \Delta$ ing~\Log{G3c}.
Kita kudu nuduhake yen uga ana \Log{G3c}-!!{proof}~$\pi'$
kanggo $\Gamma \Sequent \Delta, !A$ utawa
$!A, \Gamma \Sequent \Delta$, manut urutane, kanthi
$\pheight{\pi'} \le \pheight{\pi}$. Kita nindakake iki
kanthi induksi tumrap $n = \pheight{\pi}$.

Yen $n=0$, $\pi$ mung aksioma. Nanging yen
$\Gamma \Sequent \Delta, !A, !A$ minangka aksioma~\Log{G3c},
mula $\Gamma \Sequent \Delta, !A$ uga aksioma: ing kasus
identitas, ana formula atomik~$!C$ sing dadi !!a{element} ing $\Gamma$ lan~$\Delta$,
utawa $!A$ atomik lan dadi !!a{element} ing~$\Gamma$.
Panalaran sing padha ditrapake yen sekuen pungkasan yaiku
$!A, !A, \Gamma \Sequent \Delta$. Falsitas ing anteseden
uga tetep ana sawise kontraksi, kalebu nalika falsitas
minangka formula sing dikontraksi ing kiwa.

Yen $n>0$, $\pi$ dipungkasi nganggo aturan~$R$.
Yen $!A$ dudu formula utama ing inferensi pungkasan,
kita bisa oleh !!a{proof}~$\pi'$ kaya ing bukti
\olref{lem:G3c-invert}, kanthi ngetrapake hipotesis
induksi marang sub-!!{proof}s langsung saka~$\pi$
lan ngetrapake aturan~$R$ sing padha marang
!!{proof}(s) asile. Hipotesis induksi nyatakake yen
ana !!a{proof} kanthi !!{height}~$k<n$ kanggo
$\Pi \Sequent \Lambda, !D, !D$ utawa
$!D, !D, \Pi \Sequent \Lambda$, mula ana
!!a{proof} kanthi !!{height}~$\le k$ kanggo
$\Pi \Sequent \Lambda, !D$ utawa
$!D, \Pi \Sequent \Lambda$, manut urutane.
(Gatekna yen hipotesis induksi bisa ditrapake
sanajan konteks utawa !!{formula} sing dikontraksi
beda saka $\Gamma$, $\Delta$, utawa~$!A$!)

Upamane $R$ yaiku $\RightR{\land}$ lan $\pi$~dipungkasi nganggo
\[
\AxiomC{}
\RightLabel{$\pi_1$}
\Deduce$\Gamma \fCenter \Delta', !B, !A, !A$
\AxiomC{}
\RightLabel{$\pi_2$}
\Deduce$\Gamma \fCenter \Delta', !C, !A, !A$
\RightLabel{\RightR{\land}}
\BinaryInf$\Gamma \fCenter \Delta', !B \land !C, !A, !A$
\DisplayProof
\]
kanthi $\Delta = \Delta', !B \land !C$. Hipotesis induksi
menehi !!{proof}s $\pi_1'$ lan $\pi_2'$ kanggo
$\Gamma \fCenter \Delta', !B, !A$ lan
$\Gamma \fCenter \Delta', !C, !A$, manut urutane, kanthi
$\pheight{\pi_1'} \le \pheight{\pi_1}$ lan
$\pheight{\pi_2'} \le \pheight{\pi_2}$.
!!{proof}~$\pi'$ yaiku:
\[
\AxiomC{}
\RightLabel{$\pi_1'$}
\Deduce$\Gamma \fCenter \Delta', !B, !A$
\AxiomC{}
\RightLabel{$\pi_2'$}
\Deduce$\Gamma \fCenter \Delta', !C, !A$
\RightLabel{\RightR{\land}}
\BinaryInf$\Gamma \fCenter \Delta', !B \land !C, !A$
\DisplayProof
\]

Yen $!A$ dadi formula utama ing inferensi pungkasan,
kita ngetrapake lemma kuwalikan (\olref{lem:G3c-invert})
marang sub-!!{proof}s langsung saka~$\pi$, banjur
hipotesis induksi, lan banjur aturan~$R$ maneh.
Upamane $\pi$ mbuktekake $\Gamma \Sequent \Delta, !A, !A$,
$!A \ident (!B \land !C)$, lan formula iki dadi
formula utama ing inferensi pungkasan. Mula
$\pi$ duwe wujud iki:
\[
\AxiomC{}
\RightLabel{$\pi_1$}
\Deduce$\Gamma \fCenter \Delta, !B \land !C, !B$
\AxiomC{}
\RightLabel{$\pi_2$}
\Deduce$\Gamma \fCenter \Delta, !B \land !C, !C$
\RightLabel{\RightR{\land}}
\BinaryInf$\Gamma \fCenter \Delta, !B \land !C, !B \land !C$
\DisplayProof
\]
Trapna \olref{lem:G3c-invert} marang $\pi_1$ kanggo
oleh !!a{proof}~$\pi_1'$ kanggo
$\Gamma \Sequent \Delta, !B, !B$. (Iki uga menehi
!!a{proof} kanggo $\Gamma \Sequent \Delta, !C, !B$,
nanging bukti iki ora kita butuhake.) Semono uga,
kanthi ngetrapake \olref{lem:G3c-invert} marang~$\pi_2$,
kita oleh !!a{proof}~$\pi_2'$ kanggo
$\Gamma \Sequent \Delta, !C, !C$.
Amarga kuwalikan \RightR{\land} njaga !!{height},
$\pheight{\pi_1'} \le \pheight{\pi_1} < \pheight{\pi}$ lan
$\pheight{\pi_2'} \le \pheight{\pi_2} < \pheight{\pi}$.
Mula hipotesis induksi bisa ditrapake marang
$\pi_1'$ lan $\pi_2'$: ana !!{proof}s~$\pi_1''$
lan $\pi_2''$ kanggo $\Gamma \Sequent \Delta, !B$
lan $\Gamma \Sequent \Delta, !C$, manut urutane,
kanthi $\pheight{\pi_1''} \le \pheight{\pi_1'} \le \pheight{\pi_1}$
lan $\pheight{\pi_2''} \le \pheight{\pi_2'} \le \pheight{\pi_2}$.
!!{proof}~$\pi'$ banjur yaiku:
\[
\AxiomC{}
\RightLabel{$\pi_1''$}
\Deduce$\Gamma \fCenter \Delta, !B$
\AxiomC{}
\RightLabel{$\pi_2''$}
\Deduce$\Gamma \fCenter \Delta, !C$
\RightLabel{\RightR{\land}}
\BinaryInf$\Gamma \fCenter \Delta, !B \land !C$
\DisplayProof
\]
kanthi $\pheight{\pi'} \le \pheight{\pi}$.

Kita kudu luwih ngati-ati tumrap aturan kuantor
ing kasus nalika !!{formula} dadi formula utama.

Upamane $!A \ident \lforall[x][!B(x)]$ dadi formula
utama ing tengen. Mula $\pi$ dipungkasi nganggo
\[
\AxiomC{}
\RightLabel{$\pi_1$}
\Deduce$\Gamma \fCenter \Delta, \lforall[x][!B(x)], !B(a)$
\RightLabel{\RightR{\lforall}}
\UnaryInf$\Gamma \fCenter \Delta, \lforall[x][!B(x)], \lforall[x][!B(x)]$
\DisplayProof
\]
Miturut \olref{lem:invert-quant}, ana !!a{proof}~$\pi_1'$
kanthi !!{height} $\le \pheight{\pi_1}$ kanggo
$\Gamma \fCenter \Delta, !B(c), !B(a)$, kanggo
saben~$c$ sing ora ana ing $\Gamma \fCenter \Delta,
\lforall[x][!B(x)], !B(a)$. Kita bisa nganggep
$\pi_1'$ reguler; mula miturut
\olref[qua]{lem:subst-regular}, !!{proof}
$\Subst{\pi_1'}{a}{c}$ minangka !!a{proof} kanggo
$\Gamma \fCenter \Delta, !B(a), !B(a)$, uga kanthi
!!{height} $\le \pheight{\pi_1}$. Saiki hipotesis
induksi bisa ditrapake lan menehi !!a{proof}~$\pi_1''$
kanggo $\Gamma \fCenter \Delta, !B(a)$.
Trapna aturan \RightR{\lforall} kanggo oleh~$\pi'$:
\[
\AxiomC{}
\RightLabel{$\pi_1''$}
\Deduce$\Gamma \fCenter \Delta, !B(a)$
\RightLabel{\RightR{\lforall}}
\UnaryInf$\Gamma \fCenter \Delta, \lforall[x][!B(x)]$
\DisplayProof
\]
kanthi $\pheight{\pi'} \le \pheight{\pi}$.

Saiki upamane $!A \ident \lexists[x][!B(x)]$ lan
dadi formula utama ing tengen. Mula $\pi$ duwe wujud iki:
\[
\AxiomC{}
\RightLabel{$\pi_1$}
\Deduce$\Gamma \fCenter \Delta, \lexists[x][!B(x)], \lexists[x][!B(x)], !B(t)$
\RightLabel{\RightR{\lexists}}
\UnaryInf$\Gamma \fCenter \Delta, \lexists[x][!B(x)], \lexists[x][!B(x)]$
\DisplayProof
\]
Hipotesis induksi bisa ditrapake marang~$\pi_1$,
lan ana !!a{proof}s~$\pi_1'$ kanggo
$\Gamma \Sequent \Delta, \lexists[x][!B(x)], !B(t)$.
(Gatekna yen ing kene ora perlu lemma kuwalikan!)
!!{proof}~$\pi'$ banjur yaiku:
\[
\AxiomC{}
\RightLabel{$\pi_1'$}
\Deduce$\Gamma \fCenter \Delta,  \lexists[x][!B(x)], !B(t)$
\RightLabel{\RightR{\lexists}}
\UnaryInf$\Gamma \fCenter \Delta, \lexists[x][!B(x)]$
\DisplayProof
\]
kanthi $\pheight{\pi'} \le \pheight{\pi}$.
\end{proof}

\begin{prob}
Wenehana sketsa bukti \olref{prop:G3c-cont-adm} kanggo
\LeftR{\Contraction}. Terangna kasus nalika
$!A \ident (!B \land !C)$ lan $!A \ident (!B \lif !C)$.
\end{prob}

\begin{prob}
Wenehana sketsa bukti \olref{prop:G3c-cont-adm} kanggo
\LeftR{\Contraction} ing kasus kuantor sing isih kari,
yaiku nalika $!A \ident \lforall[x][!B(x)]$ lan
$!A \ident \lexists[x][!B(x)]$.
\end{prob}

\end{document}
